
SK 22.

Prakriti mahat tata ahamkara tasmat ca sodasaka gana 
sodasakat tasmad api pancabhyah panca bhutani

From the primordial nature evolves the Great Principle;
from this evolves the I Principle;
from this evolves the set of sixteen;
from the five of this set of sixteen evolves the five elements.

SK 23.

buddhi adhyavasaya dharma jnana viraga aisvarya
etad rupam sattvikam tasmat viparyastam tamasam

Buddhi is self-certainty.
Virtue, knowledge, dispassion, and power are its manifestations
when the sattva attribute abounds.
And the reverse of these, when the tamas attribute abounds.

SK 24.

ahamkara abhimana tasmat pravartate dvividha sarga eva
ekadasaka gana ca tanmatra pancaka

Ahamkara is self-assertion;
from that proceeds a two-fold evolution
the set of eleven and the five-fold primary elements

YS II.46

    sthira-sukham asanam
    
YS II.47

    prayatna-saithilya-ananta-samapattibhyam

YS II.48
    
    tato dvandva-anabhighata

SK 25.

vaikrita ahamkara ekadasaka sattvika pravartate
tanmatra bhutade sa tamasa taijasa ubhayam

The set of eleven abounding in sattva proceeds from the Vaikriti form of I-Principle;
the set of five primary elements proceed from the Bhutadi form of I-Principle; they are tamasa.
From the Taijasa form of I-Principle proceed both of them.

SK 26.

buddhi indriyani akhyani caksuh srotram ghrana rasana tvak 
karma indriyani ahu vak pani pada payu upastha

Organs of knowledge are called the Eye, the Ear, the Nose, the Tongue, the Skin.
Organs of action are called the Speech, the Hand, the Feet, the excretory organ and the organ of generation.

YS II.49

    tasmin sati svasa-prasvasayor gati-viccheda pranayama

YS II.50

    bahya-abhyantara-stambha-vrittir desa-kala-sankhyabhi-paridrsta dirgha-suksma

YS II.51

    bahya-abhyantara-visaya-aksepi caturtha

YS II.52
   
    tata ksiyate prakasa-avaranam

YS II.53
   
    dharanasu ca yogyata manasa

SK 27.

atra mana ubhayatmakam sankalpakam ca sadharmaya indriyam
nanatvam bahya bheda ca guna-parinama-visesa

Of these (sense organs), the Mind possesses the nature of both
(the sensory and motor organs).
It is the deliberating principle, and is also called a sense organ
since it posseses properties common to the sense organs.
Its multifariousness and also its external diversities are owing
to special modifications of the Attributes.

SK 28.

pancanam rupadisu alocanamatram isyate
pancanam vritti vacana adana viharana utsarga ca ananda

The function of the five in respect to form and the rest,
is considered to be mere observation.
Speech, manipulation, locomotion, excretion and gratification
are the functions of the other five.

YS II.54
   
    sva-visaya-asamprayoge cittasya svarupa-anukara ive indriyanam pratyahara

YS II.55

    tata parama vasyata indriyanam

SK 29.

trayasya svalaksanyam vritti sa esa asamanya bhavati
samanya karana-vritti pranadhyah-vayavah-panca

Of the three internal instruments, their own characteristics are their functions;
this is peculiar to each.
The common modification of the instruments is the five airs
such as prana and the rest.

SK 30.

drste catustayasya tu yugapat vrttih tasya kramasasca nirdista
tatha api adrste trayasya vrttih tat purvika

Of all the four, the functions are said to be simultaneous and also successive
with regard to the seen objects;
with regard to the unseen objects
the functions of the three are preceeded by that.

SK 31.

svam svam vrittim pratipadyante paraspara akuta hetukam
purushartha eva hetu na kenacit karanam karyate

Instruments enter into their respective functions being incited by mutual impulse.
The purpose of the Spirit is the sole motive (for the activities of the instruments).
By none whatsoever is an instrument made to act.

SK 32.

karanam trayodasavidham tad aharana dharana prakasakaram
tasya karyam ca dasadha aharyam dharyam prakasyam ca

Instruments are of thirteen kinds performing the functions of
seizing, sustaining and illuminating.
Its objects are of ten kinds, vis
the seized, the sustained, and the illumined.

SK 33.

antahkaranam trividham bahyam dasadha trayasya visayakhyam  
bahyam sampratakalam abhyantaram karanam trikalam

The internal instruments are three-fold.
The external are ten-fold; they are called the objects of the three (internal instruments).
The external instruments function at the present time and
the internal instruments function at all the three times.

YS III.1

    desa-bandha cittasya dharana

YS III.2

    tatra pratyaya-eka-tanata dhyanam

YS III.3

    tad evartha-matra-nirbhasa svarupa-sunyam iva samadhi

YS III.4

    trayam ekatra samyama

YS III.5

    taj-jayat prajna-aloka

YS III.6

    tasya bhumisu viniyoga

YS III.7

    trayam antar-angam purvebhya

YS III.8

    tad api bahir-angam nirbijasya

YS III.9

    vyutthana-nirodha-samskarayor abhibhava-pradur-bhavau nirodha-kshana-cittanvayo nirodha-parinama

    The Concept as such

    Manas, the understanding, first the faculty for the cognition of the general (of rules)

YS III.10

    tasya prasanta-vahita samskarat

YS III.11

    sarva-arthata-ekagratayo kshayodayau cittasya samadhi-parinama

    The logical form of all judgments consists of the objective unity of the apperception of the concepts contained therein

    Judgment

    Ahamkara, the determinative power of judgment, second the faculty for the subsumption of the particular under the general

YS III.12

    tata punashantoditau tulya-pratyayau cittasya-ekagrata-parinama

    Syllogism

    Buddhi, reason, third the faculty for the determination of the particular from the general (for the derivation from principles)

YS III.13

    etena bhutendriyesu dharma-laksanavastha-parinama vyakhyata

    Mechanism

YS III.14

    shantoditavyapadesya-dharmanupati dharmi

    Chemism

YS III.15

    kramanyatvam parinamanyatve hetu

    Teleology

YS III.16

    parinama-traya-samyamad atitanagata-jnanam


SK 34.

Tesam panca buddhi visesa-avisesa-visaya
Vak sabda-visaya-bhavati sesani tu pancavisayani

Of these, the five organs of knowledge have as their objects
both the gross and the subtle.
Speech has speech as its object;
the rest have all the five as its object.

SK 35.

yasmat buddhi santahkarana sarvam visayam avagahate
tasmat trividham karanam dvari sesani dvarani

Since buddhi along with the other internal instruments
comprehends all of the objects these three instruments are
like the warders while the rest are like the gates.

SK 36.

ete manas ahamkara pradipa-kalpa paraspara vilaksana
guna-visesa krstnam prakasya purushartham buddhau prayacchati

These (external instruments along with manas and ahamkara)
which are characteristic-wise different from one another 
and are different modifications of the attributes
and which resemble a lamp (in action) illuminating all (their respective objects)
present them to the Buddhi for the purpose of the Spirit, 

YS III.17
    sabda-artha-pratyayanam itaretaradhyasat saokaras tat-pravibhaga-samyamat sarva-bhuta-ruta-jnanam

YS III.18
    samskara-saksat-karanat purva-jati-jnanam

YS III.19
    pratyayasya para-citta-jnanam

YS III.20
    na ca tat salambanam tasya-avisayi-bhutatvat

SK 37.

yasmat buddhi purusasya upabho-gam sarvam pratisadhyati 
sa eva ca suksma pradhana-purusantaram visinasti

Because it is the Buddhi that accomplishes the experiences
with regard to all objects to the Purusha.
It is that again that discriminates the subtle difference
between the Pradhana and the Purusha.

YS III.21
    kaya-rupa-samyamat tad-grahya-sakti-stambhe caksu-prakasa-asamprayoge ‘ntardhanam

YS III.22
    etena sabdadyantardhanam uktam

SK 38.

tanmatrani avisesa tebhyo pancabhya panca bhutani
ete visesa smrta santa ghora mudha

The tanmatras are the indiscernible;
from these five proceed the five gross elements;
they are remembered as the discernible,
because, they are calm, turbulent, and deluding.

SK 39.

suksma mata-pitrja saha prabhutai tridha visesa syu 
tesam  suksma niyata mata-pitrja nivartanta

The subtle bodies, bodies born of parents, together
with gross elements are the three kinds of specific.
Of these, the subtle bodies are everlasting and
those born of the parents are perishable.

SK 40.

lingam purvotpannam asaktam niyatam mahat-adi-suksma
partantam samsarati nirupabhogam bhavai adivasitam

The mergent subtle body, produced primordially,
unconfined, constant, composed of the Tattvas beginnging with
Mahat and ending with the subtle, transmigrates, free from experience,
and tinged with dispositions.

SK 41.

yatha citram asrayam rte yatha chuya stanvadibhyo viniyoga
tadvat lingam na tisthati visesai vina nirasrayam

As a painting cannot stand without a support,
as a shadow cannot be without a stake,
similarly, the linga also cannot subsist
without a subtle body and without a support.

SK 42.

idam lingam natavat vyavastisthata purusharthahetukam
nimitta-naimittika-prasangena prakrte vibhutva yogyat

Impelled by the purpose of Purusha, this subtle body
appears in different roles, like a dramatic performer,
by means of association with instrumental causes and their effects,
through the all embracing power of Nature.
